\begin{abstract}
A formula is minimal in a logic $L$ if it belongs to $L$ and no strictly
more general formula, with respect to the substitution preorder, belongs to
$L$. Komori and Kashima asked whether intuitionistic and classical logic
are the only implicational intermediate logics axiomatizable by formulas
minimal in classical logic. Nakamura and Matsuda refuted this with an
ad hoc axiom of Gödel--Dummett logic and asked which logics are so
axiomatizable. We show that, surprisingly, all of them are: for every
classical tautology $\Phi$ we construct, in linear time, a formula $M$
minimal in classical logic such that $\IL\lplus M=\IL\lplus\Phi$, and $M$
can be taken of order at most $4$, which is optimal. The same holds for
intermediate logics in the language $\{\bot,\wedge,\vee,\to\}$. Thus
being axiomatizable by minimal formulas is not a property of a logic but
of the presentation of its axioms. Under the Curry--Howard
correspondence, a formula is minimal in $L$ exactly when it is the
principal type of every proof of it in $L$; in contrast to Hindley's
theorem that every inhabited type is the principal type of some term, our
result says that every intermediate logic is axiomatized by types that are
the principal types of all their classical inhabitants.
\end{abstract}
